\section{Experiments and Results}
\label{sec:experiments}

\subsection{Setup}
\label{sec:exp:setup}

All results in this section use the primary configuration of our
pipeline (corresponding to \texttt{adv.yaml} in the released
codebase): PaddleOCR as the detector, CoTracker3 online mode for
temporal propagation, AnyText2 as the text-editing backend, Hi-SAM
as the inpainter backend, and both the BPN and Alignment Refiner
active. Hardware is a single consumer-grade GPU. The evaluation
videos consist of real-world scenes captured with handheld or
low-rate camera motion, in natural lighting, and cover a mix of
Latin-to-Chinese and Latin-to-Spanish localization pairs.

\subsection{Overall results}
\label{sec:exp:overall}

Figure~\ref{fig:results} presents a gallery of representative
end-to-end outputs of the full pipeline across multiple input
videos and target-language pairs. Across these examples the
pipeline produces correct target glyphs, preserves perspective and
scale, and adapts to per-frame lighting and blur well enough for
the result to read as part of the original scene.
Module-level evidence for each of these properties is given in the
methodology section --- Figure~\ref{fig:lcm} for lighting,
Figure~\ref{fig:bpn:training} for blur, and
Figure~\ref{fig:refiner:training} for alignment --- and the
per-component ablations in Section~\ref{sec:exp:ablation}
quantify the contribution of each learned module.

\figplaceholder{results_overview}{\linewidth}%
  {Gallery of representative end-to-end results from the full
   pipeline across several input videos and target languages.
   \emph{Placeholder: to be filled with a curated set of good-case
   examples.}}%
  {fig:results}

\subsection{Glyph correctness: compositing vs.\ end-to-end generation}
\label{sec:exp:glyph}

Glyph-level correctness turns out to be a surprisingly demanding
requirement, and the way different systems fail on it is
instructive for understanding the design of our pipeline.

End-to-end generative video editors such as
Gen-4 Aleph~\citep{runway2025aleph} excel at photorealistic style
preservation: surface texture, lighting gradient, and motion are
reproduced in a way that our explicit-compositing pipeline does not
match. When prompted with a specific target string, however, they
frequently render unreadable or semantically incorrect glyphs,
because reliable glyph rendering conditioned on a target string is
still an open problem for current generative systems.

The same failure mode appears in a milder form inside our
pipeline's own text editor. AnyText2~\citep{tuo2025anytext2} is a
mask-based scene-text editor, and when the target string is
substantially shorter than the source, the source-derived mask
carries more space than the target glyphs need. Naively passing
this mask to AnyText2 causes it to fill the excess with
plausible-looking but semantically incorrect characters --- the same
hallucinated-glyph pattern that affects end-to-end generative
models. Our \emph{adaptive mask} path in S3
(Section~\ref{sec:method:s3}) is a deliberate engineering workaround
for this behaviour: we pre-clear the source text with the
configured background inpainter, restore a feathered centered strip
of original pixels at the target's natural width, and pass a shrunk
mask that tightly bounds the target glyphs. With this workaround,
AnyText2 reliably produces the correct target string even in
long-to-short settings, and this is what allows the compositing
path as a whole to deliver the glyph-correctness guarantee that
end-to-end generative systems currently do not.
Figure~\ref{fig:glyph:comparison} shows the three regimes side by
side.

\figplaceholder{generative_vs_adaptive}{\linewidth}%
  {Glyph correctness across approaches on a challenging scene.
   \textbf{Left:} end-to-end generative output (Runway
   Gen-4 Aleph) --- strong style preservation but incorrect target
   glyphs. \textbf{Middle:} AnyText2 applied directly with the
   source-derived mask --- gibberish characters when the target is
   shorter than the source, the same failure mode as the generative
   model in a milder form. \textbf{Right:} our pipeline with the
   adaptive-mask path enabled --- correct target glyphs, at the
   cost of coarser style preservation than the generative baseline.
   The example is chosen to illustrate the complementary strengths
   of the two paradigms (Section~\ref{sec:exp:discussion}), not as
   a head-to-head benchmark. \emph{Placeholder.}}%
  {fig:glyph:comparison}

\subsection{Module ablations}
\label{sec:exp:ablation}

To assess the contribution of the individual components we designed
or tuned, we disable them one at a time and measure three
quantitative metrics on a representative video set:
(i)~\emph{OCR readback accuracy}, defined as PaddleOCR's
string-level match rate when re-detecting the translated text in
our output; (ii)~\emph{background SSIM}~\citep{wang2004ssim}
measured outside the text alpha mask, which captures whether the
pipeline damages surrounding scene pixels; and (iii)~\emph{quad
jitter}, the mean per-frame corner displacement of the final
composite, which captures temporal stability.
Qualitative evidence for each learned component is already present
in the methodology section:
Figure~\ref{fig:bpn:training} for the BPN,
Figure~\ref{fig:refiner:training} for the Alignment Refiner,
Figure~\ref{fig:lcm} for the Lighting Correction Module,
and Figure~\ref{fig:inpaint} for the inpainter backends.
Table~\ref{tab:ablation} summarizes the quantitative contribution
of each component.

\begin{table}[H]
  \centering
  \caption{Module ablations on a representative video set. Each row
    disables a single component of the pipeline.
    OCR readback: PaddleOCR re-detection accuracy on output frames.
    Background SSIM: structural similarity in non-text regions.
    Quad jitter: mean per-frame corner displacement (lower is more
    temporally stable). \emph{All numbers below are placeholders to
    be filled with measured values.}}
  \label{tab:ablation}
  \small
  \setlength{\tabcolsep}{4pt}
  \begin{tabular}{lccc}
    \toprule
    Configuration & OCR readback $\uparrow$ & Background SSIM $\uparrow$ & Quad jitter (px) $\downarrow$ \\
    \midrule
    Full pipeline                    & XX.X & X.XX & X.X \\
    w/o Alignment Refiner            & XX.X & X.XX & X.X \\
    w/o BPN (LCM only)               & XX.X & X.XX & X.X \\
    w/o LCM (histogram match only)   & XX.X & X.XX & X.X \\
    w/o adaptive mask                & XX.X & X.XX & --- \\
    \bottomrule
  \end{tabular}
\end{table}

We expect the strongest effects to fall on temporal stability
(Alignment Refiner) and readback accuracy (adaptive mask), since
these target the most direct failure modes of a naive pipeline.
Disabling the BPN removes per-frame blur matching and leaves the
composited text uniformly sharp; we expect this to show as a
modest average SSIM gap that understates the perceptual cost,
because motion blur concentrates on a small set of high-motion
frames that averaging tends to hide.

\subsection{Analysis and insights}
\label{sec:exp:analysis}

The design choices introduced in Section~\ref{sec:method:summary}
each leave a measurable or visible trace in the results above.

\begin{itemize}[leftmargin=*]
  \item \textbf{Explicit geometric normalization reduces temporal
    jitter.} The Alignment Refiner's contribution in the
    quad-jitter column of Table~\ref{tab:ablation} makes concrete
    the cost of residual tracker error: correcting a few pixels
    of drift per frame propagates into a visibly more stable
    composite. The training-time view in
    Figure~\ref{fig:refiner:training} shows the same correction at
    the ROI level.
  \item \textbf{Stage-wise decomposition concentrates effort
    where it matters.} Editing is applied once per track at a
    carefully chosen reference frame rather than independently on
    every frame, which avoids frame-to-frame edit inconsistency
    and keeps the per-frame cost on the propagation side bounded.
    The adaptive-mask workaround in
    Section~\ref{sec:exp:glyph} is a direct consequence of this
    choice: because the editor is invoked once, we can afford to
    shape its input carefully.
  \item \textbf{Appearance-aware propagation closes the
    per-frame look-consistency gap.} LCM and BPN target
    complementary appearance axes ---
    Figure~\ref{fig:lcm} shows LCM preserving source-frame
    illumination, and Figure~\ref{fig:bpn:training} shows the BPN
    reproducing per-frame blur. Together they bring the composited
    text's per-frame appearance close enough to the surrounding
    scene that viewers do not read the edit as an overlay.
\end{itemize}

A useful side effect of this explicit decomposition is that every
intermediate representation---canonical ROIs, homography
matrices, lighting ratio maps, blur parameters---is inspectable
and individually controllable, which supports debugging and
module-level tuning in a way that end-to-end generative editors
do not.

\subsection{Discussion}
\label{sec:exp:discussion}

Compared to generative video editors such as
Gen-4~Aleph~\citep{runway2025aleph}, which prioritize global
visual realism, our pipeline prioritizes structured
controllability and glyph-level fidelity. Generative models
produce visually compelling outputs, but at the time of writing
they do not offer guarantees on target-string accuracy or precise
spatial alignment. The comparison in
Section~\ref{sec:exp:glyph} and the ablations in
Section~\ref{sec:exp:ablation} together suggest that explicitly
modelling geometry and appearance propagation is a viable route
when correct glyphs matter, while the naturalistic qualities of
generative systems remain harder to match with explicit
compositing. A reasonable reading of these observations is that
the two approaches are complementary rather than competing, and
that hybrid systems---generative backgrounds and style with
glyph-controlled text compositing---are a promising long-term
direction.

\subsection{Current limitations}
\label{sec:exp:limits}

Three limitations are worth noting explicitly. The most important
is the \emph{style-transfer gap} against end-to-end generative
baselines: by design, our pipeline produces correct target glyphs,
but the resulting texture match is not as naturalistic as Gen-4
Aleph's. We view this as the fundamental tradeoff of the
glyph-controlled paradigm. Second, \emph{temporal stability} is
imperfect even with the Alignment Refiner: homographies are
estimated per-frame and the refiner predicts per-frame
independently, which leaves small micro-oscillations in the
composited corners that are sometimes visible on long clips.
Third, the adaptive mask's character-class width heuristic is an
approximation: highly stylized fonts whose glyph widths deviate
substantially from the heuristic averages can still leave a small
residual artifact at the mask edges. All three are promising
directions for future work and do not affect the feasibility
result.
